% Options for packages loaded elsewhere
\PassOptionsToPackage{unicode}{hyperref}
\PassOptionsToPackage{hyphens}{url}
\PassOptionsToPackage{dvipsnames,svgnames,x11names}{xcolor}
\documentclass[
  11pt]{article}
\usepackage{xcolor}
\usepackage[a4paper,margin=27mm]{geometry}
\usepackage{amsmath,amssymb}
\setcounter{secnumdepth}{-\maxdimen} % remove section numbering
\usepackage{iftex}
\ifPDFTeX
  \usepackage[T1]{fontenc}
  \usepackage[utf8]{inputenc}
  \usepackage{textcomp} % provide euro and other symbols
\else % if luatex or xetex
  \usepackage{unicode-math} % this also loads fontspec
  \defaultfontfeatures{Scale=MatchLowercase}
  \defaultfontfeatures[\rmfamily]{Ligatures=TeX,Scale=1}
\fi
\usepackage{lmodern}
\ifPDFTeX\else
  % xetex/luatex font selection
\fi
% Use upquote if available, for straight quotes in verbatim environments
\IfFileExists{upquote.sty}{\usepackage{upquote}}{}
\IfFileExists{microtype.sty}{% use microtype if available
  \usepackage[]{microtype}
  \UseMicrotypeSet[protrusion]{basicmath} % disable protrusion for tt fonts
}{}
\makeatletter
\@ifundefined{KOMAClassName}{% if non-KOMA class
  \IfFileExists{parskip.sty}{%
    \usepackage{parskip}
  }{% else
    \setlength{\parindent}{0pt}
    \setlength{\parskip}{6pt plus 2pt minus 1pt}}
}{% if KOMA class
  \KOMAoptions{parskip=half}}
\makeatother
\usepackage{longtable,booktabs,array}
\usepackage{calc} % for calculating minipage widths
% Correct order of tables after \paragraph or \subparagraph
\usepackage{etoolbox}
\makeatletter
\patchcmd\longtable{\par}{\if@noskipsec\mbox{}\fi\par}{}{}
\makeatother
% Allow footnotes in longtable head/foot
\IfFileExists{footnotehyper.sty}{\usepackage{footnotehyper}}{\usepackage{footnote}}
\makesavenoteenv{longtable}
\setlength{\emergencystretch}{3em} % prevent overfull lines
\providecommand{\tightlist}{%
  \setlength{\itemsep}{0pt}\setlength{\parskip}{0pt}}
\usepackage{fancyhdr}
\usepackage{microtype}
\usepackage{titlesec}

\definecolor{MethodBlue}{HTML}{173A5E}
\definecolor{LinkBlue}{HTML}{1E5A8A}

\setlength{\parindent}{1.25em}
\setlength{\parskip}{0.35em}
\setlength{\emergencystretch}{3em}
\setlength{\headheight}{14pt}
\raggedbottom

\titleformat{\section}
  {\normalfont\LARGE\bfseries\color{MethodBlue}}
  {}{0pt}{}
\titleformat{\subsection}
  {\normalfont\Large\bfseries\color{MethodBlue}}
  {}{0pt}{}
\titlespacing*{\section}{0pt}{0pt}{1.4em}
\titlespacing*{\subsection}{0pt}{2.2em}{0.8em}

\let\originalsection\section
\renewcommand{\section}{\clearpage\originalsection}

\pagestyle{fancy}
\fancyhf{}
\fancyhead[L]{\small Expanded Welch--Satterthwaite MI}
\fancyhead[R]{\small Full derivation}
\fancyfoot[C]{\thepage}

\makeatletter
\renewcommand{\maketitle}{%
  \hypersetup{pageanchor=false}
  \begin{titlepage}
  \centering
  \vspace*{0.17\textheight}
  {\color{MethodBlue}\Huge\bfseries\@title\par}
  \vspace{2.5em}
  {\Large\normalfont\@author\par}
  \vspace{1.25em}
  {\large\normalfont\@date\par}
  \vfill
  \rule{0.72\textwidth}{0.6pt}
  \par\vspace{1em}
  \large Textbook-style mathematical derivation
  \end{titlepage}
  \hypersetup{pageanchor=true}
}
\makeatother

\newcommand{\E}{\operatorname{E}}
\newcommand{\Var}{\operatorname{Var}}
\usepackage{bookmark}
\IfFileExists{xurl.sty}{\usepackage{xurl}}{} % add URL line breaks if available
\urlstyle{same}
\hypersetup{
  pdftitle={Full Derivation of the Expanded Welch-Satterthwaite Mutual Information Test},
  pdfauthor={Daniel Zhao},
  colorlinks=true,
  linkcolor={MethodBlue},
  filecolor={Maroon},
  citecolor={Blue},
  urlcolor={LinkBlue},
  pdfcreator={LaTeX via pandoc}}

\title{Full Derivation of the Expanded Welch-Satterthwaite Mutual
Information Test}
\usepackage{etoolbox}
\makeatletter
\providecommand{\subtitle}[1]{% add subtitle to \maketitle
  \apptocmd{\@title}{\par {\large #1 \par}}{}{}
}
\makeatother
\subtitle{Two Independent Discrete Populations}
\author{Daniel Zhao}
\date{August 2026}

\begin{document}
\maketitle

{
\setcounter{tocdepth}{2}
\tableofcontents
}
\section{Purpose}\label{purpose}

This chapter derives a deterministic test for comparing the mutual
information of two independent discrete populations. The null hypothesis
is

\begin{equation}
H_0:I(P)=I(Q).
\end{equation}

The construction follows the broad logic of Welch's test:

\begin{enumerate}
\def\labelenumi{\arabic{enumi}.}
\tightlist
\item
  estimate the difference between two population quantities;
\item
  divide that difference by its estimated standard error;
\item
  account for uncertainty in the estimated standard error by using a
  Student distribution with Satterthwaite effective degrees of freedom.
\end{enumerate}

The important difference is that mutual information is a nonlinear
functional of a complete joint probability table. Its variance estimator
is therefore not an ordinary sample variance. The expanded method
derives the first-order sampling uncertainty of that complete variance
functional and uses it to calculate MI-specific component degrees of
freedom.

The derivation below is written for natural logarithms, so MI is
measured in nats.

\section{1. Statistical Setting}\label{statistical-setting}

Let

\begin{equation}
Z_1^{(P)},\ldots,Z_{n_P}^{(P)}\overset{\mathrm{iid}}{\sim}P,
\qquad
Z_1^{(Q)},\ldots,Z_{n_Q}^{(Q)}\overset{\mathrm{iid}}{\sim}Q,
\end{equation}

where every observation is a pair

\begin{equation}
Z=(X,Y)\in\{1,\ldots,r\}\times\{1,\ldots,c\}.
\end{equation}

The two samples are independent. Write the population cell probabilities
as

\begin{equation}
p_{ij}=\Pr_P(X=i,Y=j),
\qquad
q_{ij}=\Pr_Q(X=i,Y=j).
\end{equation}

For population \(P\), define the marginal probabilities

\begin{equation}
p_{i+}=\sum_{j=1}^{c}p_{ij},
\qquad
p_{+j}=\sum_{i=1}^{r}p_{ij}.
\end{equation}

The corresponding definitions apply to \(Q\).

The regular derivation assumes a fixed finite alphabet and positive
population support:

\begin{equation}
p_{ij}>0,
\qquad
q_{ij}>0
\end{equation}

for all modelled cells. This makes the logarithms and derivatives below
well-defined. Section 16 explains what changes when observed cells are
empty or the population is at independence.

\section{2. Mutual Information as a
Functional}\label{mutual-information-as-a-functional}

For a distribution \(P\), define the local-information score of cell
\((i,j)\) by

\begin{equation}
\ell_P(i,j)
=\log\left(\frac{p_{ij}}{p_{i+}p_{+j}}\right).
\end{equation}

The population mutual information is

\begin{equation}
\begin{aligned}
I(P)
&=\sum_{i=1}^{r}\sum_{j=1}^{c}
p_{ij}\log\left(\frac{p_{ij}}{p_{i+}p_{+j}}\right)\\
&=\operatorname E_P\{\ell_P(X,Y)\}.
\end{aligned}
\end{equation}

For notational convenience, write

\begin{equation}
\mu_P=I(P).
\end{equation}

The subscript on \(\ell_P\) matters. The score is not a fixed value
attached to a cell independently of the distribution. If \(P\) changes,
the joint probability, both relevant marginal probabilities, and
therefore the score all change.

\section{3. Plug-In Estimation}\label{plug-in-estimation}

Let \(N_{ij}^{(P)}\) be the observed count in cell \((i,j)\) of the
first table. Then

\begin{equation}
n_P=\sum_{i=1}^{r}\sum_{j=1}^{c}N_{ij}^{(P)},
\qquad
\widehat p_{ij}=\frac{N_{ij}^{(P)}}{n_P}.
\end{equation}

The empirical marginals are

\begin{equation}
\widehat p_{i+}=\sum_{j=1}^{c}\widehat p_{ij},
\qquad
\widehat p_{+j}=\sum_{i=1}^{r}\widehat p_{ij}.
\end{equation}

The plug-in local-information score and plug-in MI are

\begin{equation}
\widehat\ell_{ij}
=\log\left(
\frac{\widehat p_{ij}}
{\widehat p_{i+}\widehat p_{+j}}
\right)
\end{equation}

and

\begin{equation}
\widehat I(P)
=\sum_{i=1}^{r}\sum_{j=1}^{c}
\widehat p_{ij}\widehat\ell_{ij}.
\end{equation}

Equivalently, if \(\widehat P\) is the empirical distribution, then

\begin{equation}
\widehat I(P)=I(\widehat P).
\end{equation}

The same calculations produce \(\widehat I(Q)\) from the second table.

\section{4. Leading Bias Correction}\label{leading-bias-correction}

\subsection{4.1 Bias of plug-in entropy}\label{bias-of-plug-in-entropy}

For a discrete variable with \(k\) positive-probability categories, the
leading bias of the plug-in entropy estimator is

\begin{equation}
\operatorname E(\widehat H)-H
=-\frac{k-1}{2n}+O(n^{-2}).
\end{equation}

Mutual information can be written as

\begin{equation}
I(X;Y)=H(X)+H(Y)-H(X,Y).
\end{equation}

Under full support, \(X\) has \(r\) categories, \(Y\) has \(c\)
categories, and the joint variable \((X,Y)\) has \(rc\) categories.
Applying the entropy bias formula to these three terms gives

\begin{equation}
\begin{aligned}
\operatorname{Bias}(\widehat I)
&\approx
-\frac{r-1}{2n}
-\frac{c-1}{2n}
-\left(-\frac{rc-1}{2n}\right)\\
&=\frac{rc-r-c+1}{2n}\\
&=\frac{(r-1)(c-1)}{2n}.
\end{aligned}
\end{equation}

Define

\begin{equation}
d=(r-1)(c-1).
\end{equation}

\subsection{4.2 Bias-corrected MI
difference}\label{bias-corrected-mi-difference}

The leading-bias-corrected estimators are

\begin{equation}
\widehat I_{\mathrm{BC}}(P)
=\widehat I(P)-\frac{d}{2n_P}
\end{equation}

and

\begin{equation}
\widehat I_{\mathrm{BC}}(Q)
=\widehat I(Q)-\frac{d}{2n_Q}.
\end{equation}

The estimated population difference is

\begin{equation}
\widehat\Delta_{\mathrm{BC}}
=\widehat I_{\mathrm{BC}}(P)
-\widehat I_{\mathrm{BC}}(Q).
\end{equation}

Expanding this expression gives

\begin{equation}
\widehat\Delta_{\mathrm{BC}}
=\widehat I(P)-\widehat I(Q)
-\frac{d}{2n_P}+\frac{d}{2n_Q}.
\end{equation}

For fixed table dimensions and fixed sample sizes, the correction terms
are constants. They change the estimated difference but do not add
sampling variance and do not change the influence function derived next.

\section{5. Deriving the Influence Function of
MI}\label{deriving-the-influence-function-of-mi}

\subsection{5.1 Contaminate one cell}\label{contaminate-one-cell}

Fix a cell \(z=(x,y)\). Define a path of distributions

\begin{equation}
P_\varepsilon
=(1-\varepsilon)P+\varepsilon\delta_z,
\end{equation}

where \(\delta_z\) places probability one on cell \(z\). At
\(\varepsilon=0\) the distribution is \(P\). Increasing \(\varepsilon\)
moves a small amount of probability toward cell \((x,y)\) while
preserving total probability one.

For an arbitrary cell \((i,j)\),

\begin{equation}
p_{ij}(\varepsilon)
=(1-\varepsilon)p_{ij}
+\varepsilon\mathbf 1\{i=x,j=y\}.
\end{equation}

Therefore,

\begin{equation}
\left.
\frac{\mathrm d}{\mathrm d\varepsilon}p_{ij}(\varepsilon)
\right|_{\varepsilon=0}
=\mathbf 1\{i=x,j=y\}-p_{ij}.
\end{equation}

Summing over columns gives

\begin{equation}
p_{i+}(\varepsilon)
=(1-\varepsilon)p_{i+}
+\varepsilon\mathbf 1\{i=x\},
\end{equation}

so

\begin{equation}
\left.
\frac{\mathrm d}{\mathrm d\varepsilon}p_{i+}(\varepsilon)
\right|_{\varepsilon=0}
=\mathbf 1\{i=x\}-p_{i+}.
\end{equation}

Likewise,

\begin{equation}
\left.
\frac{\mathrm d}{\mathrm d\varepsilon}p_{+j}(\varepsilon)
\right|_{\varepsilon=0}
=\mathbf 1\{j=y\}-p_{+j}.
\end{equation}

\subsection{5.2 Differentiate the local-information
score}\label{differentiate-the-local-information-score}

Under \(P_\varepsilon\), the score of cell \((i,j)\) is

\begin{equation}
\ell_{P_\varepsilon}(i,j)
=\log p_{ij}(\varepsilon)
-\log p_{i+}(\varepsilon)
-\log p_{+j}(\varepsilon).
\end{equation}

Use

\begin{equation}
\frac{\mathrm d}{\mathrm d\varepsilon}\log u(\varepsilon)
=\frac{u'(\varepsilon)}{u(\varepsilon)}.
\end{equation}

The derivative of the joint-probability term at zero is

\begin{equation}
\begin{aligned}
\left.
\frac{\mathrm d}{\mathrm d\varepsilon}\log p_{ij}(\varepsilon)
\right|_{\varepsilon=0}
&=\frac{\mathbf 1\{i=x,j=y\}-p_{ij}}{p_{ij}}\\
&=\frac{\mathbf 1\{i=x,j=y\}}{p_{ij}}-1.
\end{aligned}
\end{equation}

Similarly,

\begin{equation}
\left.
\frac{\mathrm d}{\mathrm d\varepsilon}\log p_{i+}(\varepsilon)
\right|_{\varepsilon=0}
=\frac{\mathbf 1\{i=x\}}{p_{i+}}-1
\end{equation}

and

\begin{equation}
\left.
\frac{\mathrm d}{\mathrm d\varepsilon}\log p_{+j}(\varepsilon)
\right|_{\varepsilon=0}
=\frac{\mathbf 1\{j=y\}}{p_{+j}}-1.
\end{equation}

Subtracting the two marginal derivatives from the joint derivative gives

\begin{equation}
\boxed{
\dot\ell_P(i,j;z)
=
\frac{\mathbf 1\{i=x,j=y\}}{p_{ij}}
-\frac{\mathbf 1\{i=x\}}{p_{i+}}
-\frac{\mathbf 1\{j=y\}}{p_{+j}}
+1,
}
\end{equation}

where

\begin{equation}
\dot\ell_P(i,j;z)
=\left.
\frac{\mathrm d}{\mathrm d\varepsilon}
\ell_{P_\varepsilon}(i,j)
\right|_{\varepsilon=0}.
\end{equation}

\subsection{5.3 A derivative rule for distribution-dependent
expectations}\label{a-derivative-rule-for-distribution-dependent-expectations}

Suppose

\begin{equation}
F(P_\varepsilon)
=\sum_{i,j}p_{ij}(\varepsilon)h_{P_\varepsilon}(i,j),
\end{equation}

where both the probabilities and the values \(h_{P_\varepsilon}\) depend
on \(\varepsilon\). The product rule gives

\begin{equation}
\begin{aligned}
\left.\frac{\mathrm d}{\mathrm d\varepsilon}F(P_\varepsilon)
\right|_{\varepsilon=0}
={}&\sum_{i,j}
\left[\mathbf 1\{i=x,j=y\}-p_{ij}\right]h_P(i,j)\\
&+\sum_{i,j}p_{ij}\dot h_P(i,j;z).
\end{aligned}
\end{equation}

The first sum is

\begin{equation}
h_P(x,y)-\operatorname E_P\{h_P(X,Y)\}.
\end{equation}

Therefore,

\begin{equation}
\boxed{
\operatorname{IF}_{F,P}(z)
=h_P(z)-\operatorname E_P(h_P)
+\operatorname E_P\{\dot h_P(Z;z)\}.
}
\end{equation}

The last term is essential whenever the quantity being averaged changes
with the underlying distribution.

\subsection{5.4 Apply the rule to MI}\label{apply-the-rule-to-mi}

For MI, take

\begin{equation}
h_P(i,j)=\ell_P(i,j).
\end{equation}

We first evaluate the expectation of the score derivative:

\begin{equation}
\begin{aligned}
\operatorname E_P\{\dot\ell_P(Z;z)\}
={}&\sum_{i,j}p_{ij}
\left[
\frac{\mathbf 1\{i=x,j=y\}}{p_{ij}}
-\frac{\mathbf 1\{i=x\}}{p_{i+}}
-\frac{\mathbf 1\{j=y\}}{p_{+j}}
+1
\right]\\
={}&1
-\frac{\sum_j p_{xj}}{p_{x+}}
-\frac{\sum_i p_{iy}}{p_{+y}}
+\sum_{i,j}p_{ij}\\
={}&1-1-1+1\\
={}&0.
\end{aligned}
\end{equation}

The derivative rule therefore yields

\begin{equation}
\begin{aligned}
\operatorname{IF}_{I,P}(x,y)
&=\ell_P(x,y)-\operatorname E_P(\ell_P)+0\\
&=\ell_P(x,y)-\mu_P.
\end{aligned}
\end{equation}

Define

\begin{equation}
\boxed{
\psi_P(x,y)=\ell_P(x,y)-\mu_P.
}
\end{equation}

This is the influence function of mutual information.

\subsection{5.5 First-order sampling variance of plug-in
MI}\label{first-order-sampling-variance-of-plug-in-mi}

The empirical distribution satisfies a first-order functional expansion:

\begin{equation}
I(\widehat P)-I(P)
=\frac{1}{n_P}\sum_{k=1}^{n_P}\psi_P(Z_k^{(P)})
+o_p(n_P^{-1/2}).
\end{equation}

Because

\begin{equation}
\operatorname E_P\{\psi_P(Z)\}=0,
\end{equation}

the central limit theorem gives

\begin{equation}
\sqrt{n_P}\left\{\widehat I(P)-I(P)\right\}
\overset{d}{\longrightarrow}
N\{0,V(P)\},
\end{equation}

where

\begin{equation}
\begin{aligned}
V(P)
&=\operatorname{Var}_P\{\psi_P(Z)\}\\
&=\operatorname E_P\left[\{\ell_P(Z)-\mu_P\}^2\right].
\end{aligned}
\end{equation}

Consequently,

\begin{equation}
\operatorname{Var}\{\widehat I(P)\}
\approx\frac{V(P)}{n_P}.
\end{equation}

The bias correction in Section 4 is a fixed constant at a given sample
size, so it has the same first-order variance:

\begin{equation}
\operatorname{Var}\{\widehat I_{\mathrm{BC}}(P)\}
\approx\frac{V(P)}{n_P}.
\end{equation}

\section{6. The Two-Sample Standardized
Statistic}\label{the-two-sample-standardized-statistic}

Define the corresponding influence variance \(V(Q)\) for the second
population. Independence of the two samples implies

\begin{equation}
\operatorname{Var}(\widehat\Delta_{\mathrm{BC}})
\approx
\frac{V(P)}{n_P}+\frac{V(Q)}{n_Q}.
\end{equation}

To estimate \(V(P)\), define

\begin{equation}
m_{2,P}=\operatorname E_P\{\ell_P(Z)^2\}.
\end{equation}

Since \(\mu_P=\operatorname E_P(\ell_P)\),

\begin{equation}
V(P)=m_{2,P}-\mu_P^2.
\end{equation}

Its plug-in estimate is

\begin{equation}
\begin{aligned}
\widehat V_P
&=\sum_{i,j}\widehat p_{ij}
\left(\widehat\ell_{ij}-\widehat I(P)\right)^2\\
&=\sum_{i,j}\widehat p_{ij}\widehat\ell_{ij}^2
-\widehat I(P)^2.
\end{aligned}
\end{equation}

Define \(\widehat V_Q\) analogously. The two estimated contributions to
the variance of the difference are

\begin{equation}
a=\frac{\widehat V_P}{n_P},
\qquad
b=\frac{\widehat V_Q}{n_Q}.
\end{equation}

The estimated standard error is

\begin{equation}
\widehat{\operatorname{SE}}
=\sqrt{a+b}
=\sqrt{
\frac{\widehat V_P}{n_P}
+\frac{\widehat V_Q}{n_Q}
}.
\end{equation}

The standardized statistic is

\begin{equation}
\boxed{
T
=\frac{\widehat\Delta_{\mathrm{BC}}}
{\sqrt{\widehat V_P/n_P+\widehat V_Q/n_Q}}.
}
\end{equation}

If \(V(P)\) and \(V(Q)\) were known, a standard normal reference would
follow from the two independent asymptotic expansions. In practice, both
variances are estimated. Expanded Welch-Satterthwaite estimates the
uncertainty in those variance estimates rather than treating the
denominator as known.

\section{7. Deriving the Influence Function of the MI
Variance}\label{deriving-the-influence-function-of-the-mi-variance}

\subsection{7.1 The target variance
functional}\label{the-target-variance-functional}

Recall that

\begin{equation}
V(P)=m_{2,P}-\mu_P^2,
\end{equation}

where

\begin{equation}
m_{2,P}=\operatorname E_P\{\ell_P(Z)^2\},
\qquad
\mu_P=\operatorname E_P\{\ell_P(Z)\}.
\end{equation}

The goal is to differentiate \(V(P_\varepsilon)\) along the same
contamination path used in Section 5. Define

\begin{equation}
g_P(x,y)
=\left.
\frac{\mathrm d}{\mathrm d\varepsilon}V(P_\varepsilon)
\right|_{\varepsilon=0}.
\end{equation}

This \(g_P\) is the influence function of the complete MI variance
functional.

\subsection{7.2 Differentiate the second
moment}\label{differentiate-the-second-moment}

Under \(P_\varepsilon\),

\begin{equation}
m_2(P_\varepsilon)
=\sum_{i,j}p_{ij}(\varepsilon)
\ell_{P_\varepsilon}(i,j)^2.
\end{equation}

Differentiate using the product rule:

\begin{equation}
\begin{aligned}
\left.
\frac{\mathrm d}{\mathrm d\varepsilon}m_2(P_\varepsilon)
\right|_{\varepsilon=0}
={}&\sum_{i,j}
\left[\mathbf 1\{i=x,j=y\}-p_{ij}\right]\ell_P(i,j)^2\\
&+\sum_{i,j}p_{ij}
\left.
\frac{\mathrm d}{\mathrm d\varepsilon}
\ell_{P_\varepsilon}(i,j)^2
\right|_{\varepsilon=0}.
\end{aligned}
\end{equation}

For the first sum,

\begin{equation}
\sum_{i,j}
\left[\mathbf 1\{i=x,j=y\}-p_{ij}\right]\ell_P(i,j)^2
=\ell_P(x,y)^2-m_{2,P}.
\end{equation}

For the second sum, the chain rule gives

\begin{equation}
\left.
\frac{\mathrm d}{\mathrm d\varepsilon}
\ell_{P_\varepsilon}(i,j)^2
\right|_{\varepsilon=0}
=2\ell_P(i,j)\dot\ell_P(i,j;z).
\end{equation}

Therefore,

\begin{equation}
\operatorname{IF}_{m_2,P}(x,y)
=\ell_P(x,y)^2-m_{2,P}
+2\sum_{i,j}p_{ij}\ell_P(i,j)\dot\ell_P(i,j;z).
\end{equation}

\subsection{7.3 Evaluate the score-derivative
sum}\label{evaluate-the-score-derivative-sum}

Substitute the expression for \(\dot\ell_P\):

\begin{equation}
\begin{aligned}
&\sum_{i,j}p_{ij}\ell_P(i,j)\dot\ell_P(i,j;z)\\
={}&\sum_{i,j}p_{ij}\ell_P(i,j)
\left[
\frac{\mathbf 1\{i=x,j=y\}}{p_{ij}}
-\frac{\mathbf 1\{i=x\}}{p_{i+}}
-\frac{\mathbf 1\{j=y\}}{p_{+j}}
+1
\right].
\end{aligned}
\end{equation}

Evaluate the four terms separately.

The cell term is

\begin{equation}
\sum_{i,j}p_{ij}\ell_P(i,j)
\frac{\mathbf 1\{i=x,j=y\}}{p_{ij}}
=\ell_P(x,y).
\end{equation}

The row term is

\begin{equation}
\begin{aligned}
\sum_{i,j}p_{ij}\ell_P(i,j)
\frac{\mathbf 1\{i=x\}}{p_{i+}}
&=\frac{1}{p_{x+}}\sum_j p_{xj}\ell_P(x,j)\\
&=\operatorname E_P\{\ell_P(X,Y)\mid X=x\}.
\end{aligned}
\end{equation}

Define

\begin{equation}
R_P(x)
=\operatorname E_P\{\ell_P(X,Y)\mid X=x\}.
\end{equation}

The column term is

\begin{equation}
\begin{aligned}
\sum_{i,j}p_{ij}\ell_P(i,j)
\frac{\mathbf 1\{j=y\}}{p_{+j}}
&=\frac{1}{p_{+y}}\sum_i p_{iy}\ell_P(i,y)\\
&=\operatorname E_P\{\ell_P(X,Y)\mid Y=y\}.
\end{aligned}
\end{equation}

Define

\begin{equation}
C_P(y)
=\operatorname E_P\{\ell_P(X,Y)\mid Y=y\}.
\end{equation}

Finally, the constant term is

\begin{equation}
\sum_{i,j}p_{ij}\ell_P(i,j)=\mu_P.
\end{equation}

Remembering the minus signs on the row and column terms,

\begin{equation}
\boxed{
\sum_{i,j}p_{ij}\ell_P(i,j)\dot\ell_P(i,j;z)
=\ell_P(x,y)-R_P(x)-C_P(y)+\mu_P.
}
\end{equation}

It follows that

\begin{equation}
\boxed{
\begin{aligned}
\operatorname{IF}_{m_2,P}(x,y)
={}&\ell_P(x,y)^2-m_{2,P}\\
&+2\{\ell_P(x,y)-R_P(x)-C_P(y)+\mu_P\}.
\end{aligned}
}
\end{equation}

\subsection{7.4 Differentiate the squared
mean}\label{differentiate-the-squared-mean}

From Section 5,

\begin{equation}
\operatorname{IF}_{\mu,P}(x,y)
=\ell_P(x,y)-\mu_P.
\end{equation}

The ordinary chain rule for \(u^2\) gives

\begin{equation}
\left.
\frac{\mathrm d}{\mathrm d\varepsilon}
\mu(P_\varepsilon)^2
\right|_{\varepsilon=0}
=2\mu_P\operatorname{IF}_{\mu,P}(x,y).
\end{equation}

Therefore,

\begin{equation}
\boxed{
\operatorname{IF}_{\mu^2,P}(x,y)
=2\mu_P\{\ell_P(x,y)-\mu_P\}.
}
\end{equation}

\subsection{7.5 Combine the derivatives}\label{combine-the-derivatives}

Because

\begin{equation}
V(P)=m_{2,P}-\mu_P^2,
\end{equation}

linearity of differentiation gives

\begin{equation}
g_P(x,y)
=\operatorname{IF}_{m_2,P}(x,y)
-\operatorname{IF}_{\mu^2,P}(x,y).
\end{equation}

Substituting the two expressions above yields

\begin{equation}
\boxed{
\begin{aligned}
g_P(x,y)={}&
\ell_P(x,y)^2-m_{2,P}\\
&+2\{\ell_P(x,y)-R_P(x)-C_P(y)+\mu_P\}\\
&-2\mu_P\{\ell_P(x,y)-\mu_P\}.
\end{aligned}
}
\end{equation}

This is the complete variance-influence formula used by the expanded
method. No local score is treated as fixed: the middle line accounts for
changes to the joint cell, its row marginal, its column marginal, and
the overall MI.

\subsection{7.6 Verify that the influence function has mean
zero}\label{verify-that-the-influence-function-has-mean-zero}

An influence function should be centred under its defining distribution.
We can verify this directly.

For the first line,

\begin{equation}
\operatorname E_P\{\ell_P(Z)^2-m_{2,P}\}=m_{2,P}-m_{2,P}=0.
\end{equation}

For the middle line, the law of iterated expectations gives

\begin{equation}
\operatorname E_P\{R_P(X)\}
=\operatorname E_P[\operatorname E_P\{\ell_P(Z)\mid X\}]
=\mu_P
\end{equation}

and

\begin{equation}
\operatorname E_P\{C_P(Y)\}=\mu_P.
\end{equation}

Hence,

\begin{equation}
\begin{aligned}
\operatorname E_P\{\ell_P(Z)-R_P(X)-C_P(Y)+\mu_P\}
&=\mu_P-\mu_P-\mu_P+\mu_P\\
&=0.
\end{aligned}
\end{equation}

For the final line,

\begin{equation}
\operatorname E_P[-2\mu_P\{\ell_P(Z)-\mu_P\}]
=-2\mu_P(\mu_P-\mu_P)
=0.
\end{equation}

Combining the three results,

\begin{equation}
\boxed{
\operatorname E_P\{g_P(X,Y)\}=0.
}
\end{equation}

\section{8. Sampling Uncertainty of the Estimated MI
Variance}\label{sampling-uncertainty-of-the-estimated-mi-variance}

The plug-in variance estimator can be written as the empirical
functional

\begin{equation}
\widehat V_P=V(\widehat P).
\end{equation}

Because \(g_P\) is the influence function of \(V\), its first-order
expansion is

\begin{equation}
V(\widehat P)-V(P)
=\frac{1}{n_P}\sum_{k=1}^{n_P}g_P(Z_k^{(P)})
+o_p(n_P^{-1/2}).
\end{equation}

Define

\begin{equation}
\boxed{
\tau_P^2
=\operatorname{Var}_P\{g_P(X,Y)\}.
}
\end{equation}

Since \(\operatorname E_P(g_P)=0\), this can also be written as

\begin{equation}
\tau_P^2
=\operatorname E_P\{g_P(X,Y)^2\}.
\end{equation}

Taking the variance of the leading term in the expansion gives

\begin{equation}
\begin{aligned}
\operatorname{Var}\{\widehat V_P\}
&\approx
\operatorname{Var}\left\{
\frac{1}{n_P}\sum_{k=1}^{n_P}g_P(Z_k^{(P)})
\right\}\\
&=\frac{1}{n_P^2}
\sum_{k=1}^{n_P}\operatorname{Var}_P\{g_P(Z_k^{(P)})\}\\
&=\frac{1}{n_P^2}(n_P\tau_P^2)\\
&=\boxed{\frac{\tau_P^2}{n_P}}.
\end{aligned}
\end{equation}

The covariance terms vanish because observations within the sample are
independent.

This result answers the key question: \(\tau_P^2/n_P\) is the
first-order sampling variance of the estimated MI influence variance
\(\widehat V_P\).

\section{9. Estimate the Variance-Influence
Quantities}\label{estimate-the-variance-influence-quantities}

All population quantities are replaced by their empirical versions. For
the first table, calculate

\begin{equation}
\widehat\mu_P
=\sum_{i,j}\widehat p_{ij}\widehat\ell_{ij},
\end{equation}

\begin{equation}
\widehat m_{2,P}
=\sum_{i,j}\widehat p_{ij}\widehat\ell_{ij}^2,
\end{equation}

and

\begin{equation}
\widehat V_P
=\widehat m_{2,P}-\widehat\mu_P^2.
\end{equation}

For each nonempty row and column, calculate

\begin{equation}
\widehat R_P(i)
=\frac{\sum_j\widehat p_{ij}\widehat\ell_{ij}}
{\widehat p_{i+}}
\end{equation}

and

\begin{equation}
\widehat C_P(j)
=\frac{\sum_i\widehat p_{ij}\widehat\ell_{ij}}
{\widehat p_{+j}}.
\end{equation}

For every cell, evaluate

\begin{equation}
\begin{aligned}
\widehat g_P(i,j)={}&
\widehat\ell_{ij}^2-\widehat m_{2,P}\\
&+2\{\widehat\ell_{ij}-\widehat R_P(i)
-\widehat C_P(j)+\widehat\mu_P\}\\
&-2\widehat\mu_P
\{\widehat\ell_{ij}-\widehat\mu_P\}.
\end{aligned}
\end{equation}

In exact arithmetic the weighted mean of \(\widehat g_P\) is zero. The
implementation nevertheless centres it explicitly for numerical
stability:

\begin{equation}
\overline g_P
=\sum_{i,j}\widehat p_{ij}\widehat g_P(i,j).
\end{equation}

Then

\begin{equation}
\boxed{
\widehat\tau_P^2
=\sum_{i,j}\widehat p_{ij}
\{\widehat g_P(i,j)-\overline g_P\}^2.
}
\end{equation}

Repeat the complete calculation for table \(Q\) to obtain
\(\widehat V_Q\) and \(\widehat\tau_Q^2\).

\section{10. Satterthwaite Moment Matching for One Variance
Component}\label{satterthwaite-moment-matching-for-one-variance-component}

\subsection{10.1 The scaled chi-squared
approximation}\label{the-scaled-chi-squared-approximation}

A chi-squared random variable \(U\sim\chi^2_\nu\) has

\begin{equation}
\operatorname E(U)=\nu,
\qquad
\operatorname{Var}(U)=2\nu.
\end{equation}

Therefore,

\begin{equation}
Y=m\frac{U}{\nu}
\end{equation}

has

\begin{equation}
\operatorname E(Y)
=\frac{m}{\nu}\operatorname E(U)
=m
\end{equation}

and

\begin{equation}
\begin{aligned}
\operatorname{Var}(Y)
&=\frac{m^2}{\nu^2}\operatorname{Var}(U)\\
&=\frac{m^2}{\nu^2}(2\nu)\\
&=\frac{2m^2}{\nu}.
\end{aligned}
\end{equation}

Solving the last equation for \(\nu\) gives the general Satterthwaite
moment-matching rule

\begin{equation}
\boxed{
\nu=\frac{2m^2}{\operatorname{Var}(Y)}.
}
\end{equation}

\subsection{\texorpdfstring{10.2 Apply the rule to
\(\widehat V_P\)}{10.2 Apply the rule to \textbackslash widehat V\_P}}\label{apply-the-rule-to-widehat-v_p}

Approximate the positive variance estimator by

\begin{equation}
\widehat V_P
\mathrel{\dot\sim}
V(P)\frac{\chi^2_{\nu_{V,P}}}{\nu_{V,P}}.
\end{equation}

The scaled chi-squared approximation has

\begin{equation}
\operatorname E(\widehat V_P)\approx V(P)
\end{equation}

and

\begin{equation}
\operatorname{Var}(\widehat V_P)
\approx\frac{2V(P)^2}{\nu_{V,P}}.
\end{equation}

Section 8 independently derived

\begin{equation}
\operatorname{Var}(\widehat V_P)
\approx\frac{\tau_P^2}{n_P}.
\end{equation}

Equating the two approximations gives

\begin{equation}
\frac{2V(P)^2}{\nu_{V,P}}
=\frac{\tau_P^2}{n_P}.
\end{equation}

Multiply both sides by \(n_P\nu_{V,P}\):

\begin{equation}
2n_PV(P)^2
=\nu_{V,P}\tau_P^2.
\end{equation}

Divide by \(\tau_P^2\):

\begin{equation}
\boxed{
\nu_{V,P}
=\frac{2n_PV(P)^2}{\tau_P^2}.
}
\end{equation}

The empirical component degrees of freedom are therefore

\begin{equation}
\boxed{
\widehat\nu_{V,P}
=\frac{2n_P\widehat V_P^2}{\widehat\tau_P^2}.
}
\end{equation}

Likewise,

\begin{equation}
\boxed{
\widehat\nu_{V,Q}
=\frac{2n_Q\widehat V_Q^2}{\widehat\tau_Q^2}.
}
\end{equation}

\subsection{10.3 Why dividing by sample size does not change component
degrees of
freedom}\label{why-dividing-by-sample-size-does-not-change-component-degrees-of-freedom}

The denominator of \(T\) uses

\begin{equation}
A=\frac{\widehat V_P}{n_P},
\end{equation}

not \(\widehat V_P\) itself. From the scaled chi-squared approximation,

\begin{equation}
A
\mathrel{\dot\sim}
\frac{V(P)}{n_P}
\frac{\chi^2_{\nu_{V,P}}}{\nu_{V,P}}.
\end{equation}

Only the scale has changed. The chi-squared degrees of freedom remain
\(\nu_{V,P}\). The same argument applies to

\begin{equation}
B=\frac{\widehat V_Q}{n_Q}.
\end{equation}

\section{11. Combine the Two Variance
Components}\label{combine-the-two-variance-components}

The estimated squared standard error is

\begin{equation}
S^2=A+B
=\frac{\widehat V_P}{n_P}
+\frac{\widehat V_Q}{n_Q}.
\end{equation}

Let the approximate means of the two components be

\begin{equation}
a_0=\frac{V(P)}{n_P},
\qquad
b_0=\frac{V(Q)}{n_Q}.
\end{equation}

Under the component approximations,

\begin{equation}
\operatorname{Var}(A)
\approx\frac{2a_0^2}{\nu_{V,P}}
\end{equation}

and

\begin{equation}
\operatorname{Var}(B)
\approx\frac{2b_0^2}{\nu_{V,Q}}.
\end{equation}

Because the \(P\) and \(Q\) samples are independent, \(A\) and \(B\) are
independent. Thus,

\begin{equation}
\begin{aligned}
\operatorname{Var}(S^2)
&=\operatorname{Var}(A+B)\\
&=\operatorname{Var}(A)+\operatorname{Var}(B)\\
&\approx
\frac{2a_0^2}{\nu_{V,P}}
+\frac{2b_0^2}{\nu_{V,Q}}.
\end{aligned}
\end{equation}

Now approximate the sum by one scaled chi-squared variable:

\begin{equation}
S^2
\mathrel{\dot\sim}
(a_0+b_0)\frac{\chi^2_{\nu}}{\nu}.
\end{equation}

This approximation has variance

\begin{equation}
\operatorname{Var}(S^2)
\approx\frac{2(a_0+b_0)^2}{\nu}.
\end{equation}

Match the two variance expressions:

\begin{equation}
\frac{2(a_0+b_0)^2}{\nu}
=
\frac{2a_0^2}{\nu_{V,P}}
+\frac{2b_0^2}{\nu_{V,Q}}.
\end{equation}

Cancel the factor of two:

\begin{equation}
\frac{(a_0+b_0)^2}{\nu}
=
\frac{a_0^2}{\nu_{V,P}}
+\frac{b_0^2}{\nu_{V,Q}}.
\end{equation}

Solve for \(\nu\):

\begin{equation}
\nu
=\frac{(a_0+b_0)^2}
{a_0^2/\nu_{V,P}+b_0^2/\nu_{V,Q}}.
\end{equation}

Finally, replace the population components and component degrees of
freedom by their estimates:

\begin{equation}
a=\frac{\widehat V_P}{n_P},
\qquad
b=\frac{\widehat V_Q}{n_Q}.
\end{equation}

The expanded Welch-Satterthwaite degrees of freedom are

\begin{equation}
\boxed{
\widehat\nu_{\mathrm{expanded}}
=\frac{(a+b)^2}
{a^2/\widehat\nu_{V,P}
+b^2/\widehat\nu_{V,Q}}.
}
\end{equation}

\section{12. The Final Reference Distribution and
P-Value}\label{the-final-reference-distribution-and-p-value}

The expanded method compares the observed statistic

\begin{equation}
T
=\frac{\widehat\Delta_{\mathrm{BC}}}{\sqrt{a+b}}
\end{equation}

with a Student distribution having \(\widehat\nu_{\mathrm{expanded}}\)
degrees of freedom. The two-sided p-value is

\begin{equation}
\boxed{
p_{\mathrm{expanded}}
=2\left[1-F_{t_{\widehat\nu_{\mathrm{expanded}}}}
(|T|)\right],
}
\end{equation}

where \(F_{t_\nu}\) is the cumulative distribution function of a Student
random variable with \(\nu\) degrees of freedom.

At significance level \(\alpha\), reject \(H_0:I(P)=I(Q)\) when

\begin{equation}
p_{\mathrm{expanded}}<\alpha.
\end{equation}

Equivalently, reject when

\begin{equation}
|T|>
t_{1-\alpha/2,\widehat\nu_{\mathrm{expanded}}}.
\end{equation}

\section{13. Complete Calculation from Two Count
Tables}\label{complete-calculation-from-two-count-tables}

For clarity, the entire procedure is collected here without omitting any
intermediate quantity.

\subsection{Step 1: Convert counts to
probabilities}\label{step-1-convert-counts-to-probabilities}

For each group \(G\in\{P,Q\}\),

\begin{equation}
\widehat p^{(G)}_{ij}=\frac{N^{(G)}_{ij}}{n_G}.
\end{equation}

\subsection{Step 2: Calculate row and column
marginals}\label{step-2-calculate-row-and-column-marginals}

\begin{equation}
\widehat p^{(G)}_{i+}=\sum_j\widehat p^{(G)}_{ij},
\qquad
\widehat p^{(G)}_{+j}=\sum_i\widehat p^{(G)}_{ij}.
\end{equation}

\subsection{Step 3: Calculate local-information
scores}\label{step-3-calculate-local-information-scores}

\begin{equation}
\widehat\ell^{(G)}_{ij}
=\log\left(
\frac{\widehat p^{(G)}_{ij}}
{\widehat p^{(G)}_{i+}\widehat p^{(G)}_{+j}}
\right).
\end{equation}

\subsection{Step 4: Calculate plug-in
MI}\label{step-4-calculate-plug-in-mi}

\begin{equation}
\widehat I(G)
=\sum_{i,j}\widehat p^{(G)}_{ij}
\widehat\ell^{(G)}_{ij}.
\end{equation}

\subsection{Step 5: Correct the leading
bias}\label{step-5-correct-the-leading-bias}

\begin{equation}
\widehat I_{\mathrm{BC}}(G)
=\widehat I(G)-\frac{(r-1)(c-1)}{2n_G}.
\end{equation}

\subsection{Step 6: Form the estimated
difference}\label{step-6-form-the-estimated-difference}

\begin{equation}
\widehat\Delta_{\mathrm{BC}}
=\widehat I_{\mathrm{BC}}(P)
-\widehat I_{\mathrm{BC}}(Q).
\end{equation}

\subsection{Step 7: Calculate the first two score
moments}\label{step-7-calculate-the-first-two-score-moments}

For each group,

\begin{equation}
\widehat\mu_G
=\sum_{i,j}\widehat p^{(G)}_{ij}
\widehat\ell^{(G)}_{ij},
\end{equation}

\begin{equation}
\widehat m_{2,G}
=\sum_{i,j}\widehat p^{(G)}_{ij}
\left(\widehat\ell^{(G)}_{ij}\right)^2,
\end{equation}

and

\begin{equation}
\widehat V_G
=\widehat m_{2,G}-\widehat\mu_G^2.
\end{equation}

\subsection{Step 8: Calculate conditional score
means}\label{step-8-calculate-conditional-score-means}

\begin{equation}
\widehat R_G(i)
=\frac{\sum_j\widehat p^{(G)}_{ij}
\widehat\ell^{(G)}_{ij}}
{\widehat p^{(G)}_{i+}},
\end{equation}

\begin{equation}
\widehat C_G(j)
=\frac{\sum_i\widehat p^{(G)}_{ij}
\widehat\ell^{(G)}_{ij}}
{\widehat p^{(G)}_{+j}}.
\end{equation}

\subsection{Step 9: Calculate the variance influence in each
cell}\label{step-9-calculate-the-variance-influence-in-each-cell}

\begin{equation}
\begin{aligned}
\widehat g_G(i,j)={}&
\left(\widehat\ell^{(G)}_{ij}\right)^2
-\widehat m_{2,G}\\
&+2\left\{
\widehat\ell^{(G)}_{ij}
-\widehat R_G(i)-\widehat C_G(j)+\widehat\mu_G
\right\}\\
&-2\widehat\mu_G
\left\{\widehat\ell^{(G)}_{ij}-\widehat\mu_G\right\}.
\end{aligned}
\end{equation}

\subsection{Step 10: Calculate variance-influence
variability}\label{step-10-calculate-variance-influence-variability}

\begin{equation}
\overline g_G
=\sum_{i,j}\widehat p^{(G)}_{ij}\widehat g_G(i,j),
\end{equation}

\begin{equation}
\widehat\tau_G^2
=\sum_{i,j}\widehat p^{(G)}_{ij}
\left\{\widehat g_G(i,j)-\overline g_G\right\}^2.
\end{equation}

\subsection{Step 11: Calculate component degrees of
freedom}\label{step-11-calculate-component-degrees-of-freedom}

\begin{equation}
\widehat\nu_{V,G}
=\frac{2n_G\widehat V_G^2}{\widehat\tau_G^2}.
\end{equation}

\subsection{Step 12: Calculate the two standard-error
components}\label{step-12-calculate-the-two-standard-error-components}

\begin{equation}
a=\frac{\widehat V_P}{n_P},
\qquad
b=\frac{\widehat V_Q}{n_Q}.
\end{equation}

\subsection{Step 13: Calculate the standard error and
statistic}\label{step-13-calculate-the-standard-error-and-statistic}

\begin{equation}
\widehat{\operatorname{SE}}=\sqrt{a+b},
\end{equation}

\begin{equation}
T=\frac{\widehat\Delta_{\mathrm{BC}}}
{\widehat{\operatorname{SE}}}.
\end{equation}

\subsection{Step 14: Combine the component degrees of
freedom}\label{step-14-combine-the-component-degrees-of-freedom}

\begin{equation}
\widehat\nu_{\mathrm{expanded}}
=\frac{(a+b)^2}
{a^2/\widehat\nu_{V,P}
+b^2/\widehat\nu_{V,Q}}.
\end{equation}

\subsection{Step 15: Calculate the two-sided
p-value}\label{step-15-calculate-the-two-sided-p-value}

\begin{equation}
p_{\mathrm{expanded}}
=2\left[1-F_{t_{\widehat\nu_{\mathrm{expanded}}}}(|T|)\right].
\end{equation}

\section{14. Interpretation of the Effective Degrees of
Freedom}\label{interpretation-of-the-effective-degrees-of-freedom}

The component formula is

\begin{equation}
\widehat\nu_{V,P}
=\frac{2n_P\widehat V_P^2}{\widehat\tau_P^2}.
\end{equation}

Its numerator contains sample size and the squared magnitude of the MI
variance. Its denominator measures how strongly that estimated variance
changes when probability mass moves among cells.

If the cell sensitivities \(\widehat g_P(i,j)\) are similar, then
\(\widehat\tau_P^2\) is small. The estimated variance is relatively
stable, the component degrees of freedom are large, and the Student
reference is close to a normal reference.

If a small number of cells have unusually large sensitivities, then
\(\widehat\tau_P^2\) is large. The estimated variance is less stable,
the component degrees of freedom are smaller, and the Student reference
develops heavier tails. For a fixed observed \(|T|\), heavier tails
produce a larger, more cautious p-value.

This mechanism is table-dependent. Unlike the simple assignment
\(\nu_{V,P}=n_P-1\), it can respond to skewness, sparse empirical
support, and unequal influence among cells.

\section{15. Relationship to Ordinary
Welch-Satterthwaite}\label{relationship-to-ordinary-welch-satterthwaite}

An ordinary Welch test uses sample variances whose exact scaled
chi-squared distributions are available under normal sampling. A
conventional component with sample size \(n_P\) is assigned \(n_P-1\)
degrees of freedom, leading to

\begin{equation}
\widehat\nu_{\mathrm{simple}}
=\frac{(a+b)^2}
{a^2/(n_P-1)+b^2/(n_Q-1)}.
\end{equation}

Expanded Welch keeps the same Satterthwaite combination but replaces the
ordinary component degrees of freedom with quantities derived from the
MI variance functional:

\begin{equation}
n_P-1
\quad\longrightarrow\quad
\widehat\nu_{V,P}
=\frac{2n_P\widehat V_P^2}{\widehat\tau_P^2},
\end{equation}

and similarly for \(Q\).

Thus, the expansion is not a different standard error and not a second
bias correction. It is an MI-specific calculation of how uncertain each
estimated variance component is.

\section{16. Regularity Conditions and Limits of the
Derivation}\label{regularity-conditions-and-limits-of-the-derivation}

\subsection{16.1 Fixed finite alphabet}\label{fixed-finite-alphabet}

The asymptotic expansions assume that \(r\) and \(c\) remain fixed as
the sample sizes increase. The derivation does not establish validity
when the alphabet grows with sample size.

\subsection{16.2 Positive population
support}\label{positive-population-support}

The derivatives of \(\log p_{ij}\) require positive population
probabilities. The mathematical derivation therefore applies to cells in
the fixed positive support of the population. Structural-zero models
require a separate support definition and a corresponding bias
dimension.

In computation, an observed zero-count cell contributes zero to
probability-weighted sums under the convention \(0\log 0=0\). This keeps
the plug-in calculation finite, but it does not remove the underlying
smoothness assumption. Very sparse tables remain a finite-sample stress
regime rather than a setting in which the approximation becomes exact.

\subsection{16.3 Nondegenerate first-order MI
variance}\label{nondegenerate-first-order-mi-variance}

At exact independence,

\begin{equation}
p_{ij}=p_{i+}p_{+j},
\end{equation}

so

\begin{equation}
\ell_P(i,j)=0,
\qquad
\psi_P(i,j)=0,
\qquad
V(P)=0.
\end{equation}

The first-order normal approximation then degenerates. The statistic in
this chapter requires a positive combined first-order variance and is
intended for regular differential-MI comparisons away from this
degeneracy. A test whose null is independence requires second-order
theory and is not supplied by this derivation.

\subsection{16.4 Independent samples}\label{independent-samples}

The variance addition and the Satterthwaite combination use independence
of the \(P\) and \(Q\) samples. Paired, clustered, repeated-measures, or
otherwise dependent samples require covariance terms and a different
derivation.

\subsection{16.5 Fixed table dimensions in the bias
correction}\label{fixed-table-dimensions-in-the-bias-correction}

The implemented leading correction uses

\begin{equation}
d=(r-1)(c-1)
\end{equation}

from the configured table dimensions. If categories are selected or
removed after observing the data, the nominal bias calculation and the
sampling analysis can change.

\subsection{16.6 The Student reference is an
approximation}\label{the-student-reference-is-an-approximation}

The derivation establishes first-order moments for the variance
estimator and uses Satterthwaite moment matching. It does not prove that
\(\widehat V_P\) is exactly scaled chi-squared. It also does not make
the MI numerator exactly independent of its estimated denominator.

Consequently,

\begin{equation}
T\not\equiv t_{\widehat\nu_{\mathrm{expanded}}}
\end{equation}

as an exact finite-sample identity. The Student distribution is a
calibrated working reference motivated by the derived variance
uncertainty. Its finite-sample accuracy must be established empirically.

\subsection{16.7 Numerical validity
conditions}\label{numerical-validity-conditions}

The calculation requires finite positive values for

\begin{equation}
\widehat V_P+\widehat V_Q,
\qquad
\widehat\tau_P^2,
\qquad
\widehat\tau_Q^2,
\end{equation}

and for the resulting component and combined degrees of freedom. If
these conditions fail, the implementation reports the expanded result as
invalid rather than manufacturing a p-value.

\section{17. Computational Complexity}\label{computational-complexity}

Each quantity is obtained by a fixed number of operations over an
\(r\times c\) table:

\begin{itemize}
\tightlist
\item
  cell probabilities and local scores require \(O(rc)\) work;
\item
  row and column reductions require \(O(rc)\) work;
\item
  \(\widehat V\), \(\widehat g\), and \(\widehat\tau^2\) require
  \(O(rc)\) work;
\item
  combining the two components requires constant work.
\end{itemize}

Therefore, the complete two-table method has

\begin{equation}
\boxed{
\text{time complexity }O(rc),
\qquad
\text{memory complexity }O(rc).
}
\end{equation}

No permutations, bootstrap samples, or Monte Carlo tables are required.

\section{18. Correspondence with the
Implementation}\label{correspondence-with-the-implementation}

The implementation is in
\href{../src/welch_differential_mi/welch.py}{\texttt{src/welch\_differential\_mi/welch.py}}.
Its main quantities correspond to the derivation as follows.

{\def\LTcaptype{} % do not increment counter
\begin{longtable}[]{@{}
  >{\raggedright\arraybackslash}p{(\linewidth - 2\tabcolsep) * \real{0.5000}}
  >{\raggedright\arraybackslash}p{(\linewidth - 2\tabcolsep) * \real{0.5000}}@{}}
\toprule\noalign{}
\begin{minipage}[b]{\linewidth}\raggedright
Mathematical quantity
\end{minipage} & \begin{minipage}[b]{\linewidth}\raggedright
Implementation name
\end{minipage} \\
\midrule\noalign{}
\endhead
\bottomrule\noalign{}
\endlastfoot
\(\widehat I\) & \texttt{plugin\_mi(...)} \\
\(d\) & \texttt{mi\_df} \\
\(\widehat\Delta_{\mathrm{BC}}\) & \texttt{delta} \\
\(\widehat V_P,\widehat V_Q\) & \texttt{variance\_p},
\texttt{variance\_q} \\
\(a,b\) & \texttt{component\_p}, \texttt{component\_q} \\
\(T\) & \texttt{statistic} \\
\(\widehat R(i)\) & \texttt{row\_score\_mean} \\
\(\widehat C(j)\) & \texttt{column\_score\_mean} \\
\(\widehat g(i,j)\) & \texttt{variance\_influence} \\
\(\widehat\tau^2\) & \texttt{influence\_variance} inside
\texttt{\_variance\_influence\_component\_df} \\
\(\widehat\nu_{V,P},\widehat\nu_{V,Q}\) & \texttt{expanded\_df\_p},
\texttt{expanded\_df\_q} \\
\(\widehat\nu_{\mathrm{expanded}}\) & \texttt{expanded\_df} \\
\(p_{\mathrm{expanded}}\) & \texttt{expanded\_p} \\
\end{longtable}
}

\section{19. Derivation in One Chain}\label{derivation-in-one-chain}

The complete mathematical chain is

\begin{equation}
P
\longrightarrow
\ell_P
\longrightarrow
I(P)=\mu_P
\longrightarrow
\psi_P=\ell_P-\mu_P
\longrightarrow
V(P)=\operatorname{Var}_P(\psi_P)
\end{equation}

followed by

\begin{equation}
V(P)
\longrightarrow
g_P=\operatorname{IF}_{V,P}
\longrightarrow
\tau_P^2=\operatorname{Var}_P(g_P)
\longrightarrow
\nu_{V,P}=\frac{2n_PV(P)^2}{\tau_P^2}.
\end{equation}

After repeating the chain for \(Q\),

\begin{equation}
(\nu_{V,P},\nu_{V,Q})
\longrightarrow
\nu_{\mathrm{expanded}}
\longrightarrow
t_{\nu_{\mathrm{expanded}}}
\longrightarrow
p_{\mathrm{expanded}}.
\end{equation}

The method therefore uses one influence function to estimate the
sampling variance of MI and a second influence function to estimate the
sampling uncertainty of that variance. Satterthwaite moment matching
converts the second uncertainty calculation into the effective degrees
of freedom used to interpret the standardized MI difference.

\section{References}\label{references}

\begin{itemize}
\tightlist
\item
  Hutcheson, K. (1970). \emph{A Test for Comparing Diversities Based on
  the Shannon Formula}. Journal of Theoretical Biology, 29, 151-154.
  \url{https://doi.org/10.1016/0022-5193(70)90124-4}
\item
  Satterthwaite, F. E. (1946). \emph{An Approximate Distribution of
  Estimates of Variance Components}. Biometrics Bulletin, 2, 110-114.
  \url{https://doi.org/10.2307/3002019}
\item
  Welch, B. L. (1947). \emph{The Generalization of Student's Problem
  When Several Different Population Variances Are Involved}. Biometrika,
  34, 28-35. \url{https://doi.org/10.1093/biomet/34.1-2.28}
\end{itemize}

\end{document}
